\begin{proposition}[7.1.2]
\label{IV.2.s7.1.2-fr}
Sean $A$ un anillo local noetheriano y $\mathfrak p_i$
\textup{(}$1\leqslant i\leqslant
n$\textup{)} sus ideales primos minimales.
Para que $A$ sea formalmente equidimensional es necesario
y suficiente que los $A/\mathfrak p_i$ lo sean y que $A$
sea equidimensional \textup{(es decir, que los $A/\mathfrak p_i$
tengan la misma dimensión)}.
\end{proposition}

\begin{proof}
En efecto, esto se deduce de que, para todo ideal primo
$\mathfrak p'$ de $\widehat A$, $\mathfrak p'\cap A$
contiene uno de los $\mathfrak p_i$; por tanto,
$\mathfrak p'$ contiene uno de los
$\mathfrak p_i\widehat A$, y
$\widehat A/\mathfrak p_i\widehat A=(A/\mathfrak p_i)\otimes_A\widehat A$
es el completado del anillo local $A/\mathfrak p_i$
(\textbf 0$_{\mathrm I}$, 7.3.3), y por tanto tiene la misma
dimensión. Así, toda cadena maximal de ideales primos de
$\widehat A$ se identifica canónicamente con una cadena
maximal de ideales primos de uno de los
$\widehat A/\mathfrak p_i\widehat A$, y recíprocamente,
de donde se sigue inmediatamente la conclusión.
\end{proof}
